%%%%

%%%   Самарский государственный технический университет

%%%   Кафедра прикладной математики и информатики

%%%

%%%   Преамбула для статьи в журнал "Вестник Самарского государственного

%%%   технического университета. Серия: «Физико-математические науки»"

%%%   Ориентирована под MikTEx 2.4 for Win32

%%%

%%%   Хотя сам журнал собирается под GNU/Linux на дистрибутиве TeXLive



\documentclass[11pt,a4paper,twoside]{article}



\usepackage[cp1251]{inputenc}

\usepackage[T2A]{fontenc}

\usepackage[english,russian]{babel}

\usepackage{samgtubib}

\usepackage[dvips]{graphicx}

\usepackage[multiple]{footmisc}



\usepackage{samgtu}

\renewcommand{\VestnikYear}{2009} % Год издания Вестника

\renewcommand{\VestnikNumber}{2\,(19)} % Номер Вестника

\renewcommand{\VestnikMonth}{Ноябрь} % Месяц выхода Вестника

\renewcommand{\VestnikData}{01.11.09} % Дата подписания Вестника в печать

\renewcommand{\VestnikUPL}{26,64} % Усл. печ. л.

\renewcommand{\VestnikUIL}{25,73} % Уч.-изд. л.

\renewcommand{\VestnikRegNumber}{479/09} % Регистрационный номер

\renewcommand{\VestnikOrder}{994} % Номер заказа







%

%\begin{document}





\renewcommand{\firstpage}{185}

\renewcommand{\lastpage}{189}



\udk{517.956.6}

\msc{35M13}



\title{Об одной краевой задаче  для~уравнения смешанного типа с~нелокальными начальными условиями в~прямоугольнике}{Об одной краевой задаче  для~уравнения смешанного типа  \ldots}



\author{С.~В.~Кириченко}{С.~В.~Кириченко}





\institute{\samgups.}



\email{E-mail}{svkirichenko@mail.ru}





\otherinf{\fullauthor{Светлана Викторовна Кириченко} \position{старший преподаватель} \dept{каф. математики} }



\keywords{уравнение смешанного типа, нелокальные условия, обобщённое решение}



\workabstract{Рассмотрена краевая задача с нелокальным начальным условием для  уравнения смешанного типа. Основным результатом является доказательство эквивалентности поставленной нелокальной задачи

и краевой задачи с классическими начальными условиями для нагруженного уравнения.

Установленная эквивалентность позволила доказать единственность решения поставленной задачи

и существование обобщённого решения при дополнительных ограничениях на входные данные.}





\engtitle{On a boundary value problem for mixed type equation with nonlocal initial conditions in the rectangle}



\engauthor{S.~V.~Kirichenko}{S.~V.~K\,i\,r\,i\,c\,h\,e\,n\,k\,o}



\enginstitute{\engsamgups.}



\otherenginf{\fullauthor{Svetlana V. Kirichenko}

\position{Senior Teacher} \dept{Dept. of Mathematics}}



\engabstract{The boundary value problem for mixed type equation with nonlocal initial conditions in integral form is considered. The main result states that the nonlocal problem is equivalent to the classical boundary value problem for a loaded equation. This fact helps to prove the uniqueness and, under extra restrictions, the existence of a generalized solution of the problem.}



\engkeywords{mixed type equation, nonlocal conditions, generalized solution}



\date{24/VII/2013}

\revisiondate{09/VIII/2013}



\maketitle



\small



\Section[n]{Введение} В настоящее время задачи с нелокальными

условиями для уравнений в частных производных активно изучаются.

Интерес к ним вызван необходимостью обобщения классических задач

математической физики в связи с математическим моделированием ряда

физических процессов, изучаемых современным естествознанием

\cite{kir:sam}. В тех случаях, когда граница области протекания

физического процесса недоступна для непосредственных измерений,

дополнительной информацией, достаточной для однозначной

разрешимости соответствующей математической задачи, могут служить

нелокальные условия в интегральной форме. Исследованию

разрешимости нелокальной задачи с интегральными условиями и

посвящена данная статья.



Значительная часть публикаций, начиная с работы Дж.~Р.~Кэннона

\cite{kir:cannon}, посвященных задачам с интегральными условиями,

содержит исследования параболических уравнений. Важные результаты

исследования нелокальных задач для эллиптических уравнений

получены А.~К.~Гущиным, В.~П.~Михайловым~\cite{kir:Gu},

А.~Л.~Скубачевским~\cite{kir:skub}. Задачи с интегральными

условиями для гиперболических уравнений рассмотрены в статьях

[\citen{kir:G-A}--\citen{kir:Pu_1}], а также в работах,

представленных в прилагающихся  к ним списках литературы. Нелокальным задачам для

уравнений смешанного типа посвящено весьма небольшое количество

работ, причем в них рассмотрены лишь модельные уравнения. Отметим

здесь работы К.~Б.~Сабитова (см., например, \cite{kir:sabit}) и

его учеников.



В предлагаемой статье рассмотрена задача с нелокальным по времени интегральным условием для уравнения смешанного типа.



\smallskip



\Section[n]{Постановка задачи} В конечной области $Q_T=(0,l)\times

(0,T)$ рассмотрим уравнение

\begin{equation}

\label{kir:eq1} Lu\equiv

K(x,t)u_{tt}-(a(x,t)u_x)_x+b(x,t)u_t+c(x,t)u=f(x,t).

\end{equation}



\smallskip



\begin{task}[1]Найти в области $Q_T$ решение уравнения \eqref{kir:eq1}$,$ удовлетворяющее условиям

\begin{gather}

\label{kir:2} u_x(0,t)=u_x(l,t)=0,\\

\label{kir:4} u_t(x, 0)=0,\\

\label{kir:4} u(x, 0)+ \int _0^T H(t)u(x,t)dt=0.

\end{gather}

В условии \eqref{kir:4} $H(t)$ задана в $[0, T]$ и такова$,$ что $h=\|H\|_{L_2(0, T)}<1.$ 

Коэффициент $K(x,t)$ может обращаться в

нуль как на границе области $Q_T,$ так и во внутренних её точках$.$

Мы не делаем предположений о том, где и как внутри области функция

$K(x,t)$ меняет знак$.$

\end{task}



\smallskip



\Section[n]{Эквивалентность} Покажем, что задача \eqref{kir:eq1}--\eqref{kir:4} эквивалентна задаче с классическими

начальными данными для нагруженного уравнения.



Введём оператор $B$ формулой

$$

Bu\equiv u(x,t)+ \int _0^TH(t)u(x,t)dt

$$

и будем обозначать $v(x,t)=Bu$.



Пусть $u(x,t)$ "--- решение задачи

\eqref{kir:eq1}--\eqref{kir:4}. Тогда, как нетрудно видеть,

функция $v=Bu$ удовлетворяет условиям

$$

v_x(0,t)=v_x(l,t)=0, \quad v(x, 0)=0, \quad  v_t(x, 0)=0

$$

и уравнению

$$

Lv-L\biggl( \int _0^TH(t)u(x,t)dt \biggr)=f(x,t).

$$

Преобразуем последнее слагаемое левой части этого уравнения:

$$

L \biggl( \int _0^TH(t)u(x,t)dt \biggr)=- \int _0^T

H(\tau)(au_x)_xd\tau +c(x,t) \int _0^TH(\tau)u(x,\tau)d\tau

.

$$

Так как по предположению $u(x,t)$ удовлетворяет уравнению

\eqref{kir:eq1}, 

$$

(au_x)_x=u_{tt}+bu_t+cu-f(x,t).

$$

Это представление дает нам возможность сделать полезные

преобразования:

\begin{multline*}

 \int _0^TH(\tau)(au_x)_xd\tau= \int _0^TH(\tau)K(x,\tau)u_{\tau

\tau}d\tau+\\

+ \int _0^TH(\tau)b(x,\tau)u_{\tau}+ \int _0^TH(\tau)c(x,\tau)u(x,\tau)d\tau- \int _0^THfd\tau.

\end{multline*}

Первые два слагаемые правой части последнего равенства

проинтегрируем от $0$ до $T$ и, учитывая условия

\eqref{kir:2}--\eqref{kir:4}, получим

\begin{multline*}

 \int _0^TH(\tau)(au_x)_xd\tau= \int _0^T(HK)_{\tau\tau}ud\tau-

H(0)K_t(x, 0) \int _0^THud\tau-\\

- \int _0^T(Hb)_{\tau}ud\tau+H(0)b(x, 0) \int _0^THc(x,\tau)ud\tau- \int _0^THfd\tau.

\end{multline*}

Тогда, обозначив

\begin{gather*}

P(x,t,\tau)=[c(x,t){-}c(x,\tau)]H(\tau){-}(HK)_{\tau\tau}{+}(Hb)_{\tau}{+}H(0)[K(x, 0){+}b(x, 0)]H(\tau),\\

F(x,t)=f(x,t)+ \int _0^TH(t)u(x,t)dt,

\end{gather*}

приходим к уравнению

\begin{equation}

\label{kir:5}

K(x,t)v_{tt}-(av_x)_x+bv_t+c(x,t)v+ \int _0^TP(x,t,\tau)u(x,\tau)d\tau=F(x,t).

\end{equation}

Выше мы отметили, что функция $v(x,t)$ удовлетворяет классическим условиям.



Итак, показано, что если функция $u(x,t)$ является решением задачи

\eqref{kir:eq1}--\eqref{kir:4}, то функция $v(x,t)$ "--- решение

уравнения \eqref{kir:5}, удовлетворяющее условиям

\begin{equation}

\label{kir:6} v_x(0,t)=v_x(l,t)=0, \quad v(x, 0)=0, \quad

v_t(x, 0)=0,

\end{equation}

а функции $u$, $v$ связаны соотношением

\begin{equation}

 \label{kir:7}

u(x,t)+ \int _0^TH(t)u(x,t)dt=v(x,t).

\end{equation}



Пусть теперь $v(x,t)$ "--- решение задачи

\eqref{kir:5}--\eqref{kir:7}. Заметим, что уравнение \eqref{kir:5}

можно записать в виде

$$

LBu-L(Bu-u)=Bf-f.

$$

Из этого равенства моментально следует, что $Lu=f$. Из условий

\eqref{kir:6} и соотношения \eqref{kir:7} следует выполнение

условий \eqref{kir:2}--\eqref{kir:4}. Эквивалентность доказана.



Таким образом, для обоснования разрешимости задачи 1 может быть рассмотрена эквивалентная ей



\smallskip



\begin{task}[2]Найти в области $Q_T$ решение уравнения

\eqref{kir:5}$,$ удовлетворяющее условиям \eqref{kir:6}$,$ если

функции $u(x,t),$ $v(x,t)$ связаны равенством \eqref{kir:7}$.$

\end{task}



\smallskip



Введём обозначения:

$$

p=

\max\limits_{[0, l]}\biggl( \int _0^T \int _0^TP^2(x,t,\tau)dtd\tau

\biggr)^{{1}/{2}}, \quad h=\|H\|_{L_2(0,T)}.

$$



\smallskip



\begin{theorem}[1]Пусть

$  c\in C(\bar Q_T), $ $  c_t\in C(\bar{Q}_T), $ $ c(x,T)\geq 0,$ $h<1,$ $2b-K_t>1,$

и либо

$$

 -a_t\geq 0, \quad p+c_t(1-\sqrt{hT})^2<0,

$$

либо

$$ -a_t>0, \quad p+c_t(1-\sqrt{hT})^2\leq 0.

$$

Тогда существует не более одного решения задачи~$2.$

\end{theorem}



\smallskip



\begin{proof} Так как уравнения \eqref{kir:5} и \eqref{kir:7} линейны, 

достаточно показать, что соответствующая однородная задача имеет

только тривиальное решение.



Умножим равенство \eqref{kir:5} с $F(x,t)=0$ на $v_t$ и

проинтегрируем по $Q_T$. После стандартных преобразований получим

\begin{multline}

 \label{kir:8}

 \int _0^T \int _0^l[(2b-K_t)v_t^2-a_tv_x^2-c_tv^2]dxdt+ \int _0^la(x,T)v_x^2(x,T)dx+\\

+ \int _0^lc(x,T)v^2(x,T)dx=

-2 \int _0^T \int _0^lv_t(x,t) \int _0^TP(x,t,\tau)u(x,\tau)d\tau

dxdt.

\end{multline}

Для оценки правой части \eqref{kir:8} получим предварительно

неравенство. Умножим обе части равенства

$$

u(x,t)=v(x,t)- \int _0^TH(t)u(x,t)dt

$$

на $u(x,t)$ скалярно. Из полученного равенства вытекает неравенство

\begin{equation}

 \label{kir:9}

\|u\|_{L_2(Q_T)}\leq {(1-\sqrt{hT})^{-1}}\|v\|_{L_2(Q_T)}.

\end{equation}



Теперь мы можем получить оценку правой части \eqref{kir:8}:

\begin{multline*}

2\biggl| \int _0^T \int _0^lv_t(x,t) \int _0^TP(x,t,\tau)u(x,\tau)d\tau

dxdt\biggr|\leq  \int _0^T \int _0^lv_t^2dxdt+\\

+ \int _0^T \int _0^l\biggl( \int _0^TP(x,t,\tau)u(x,\tau)d\tau

\biggr)^2dxdt \leq

  \int _0^T \int _0^lv_t^2dxdt+p \int _0^T \int _0^lu^2dxdt.

\end{multline*}



В силу неравенства \eqref{kir:9} окончательно получим

\begin{multline}

 \label{kir:10}

2\biggl| \int _0^T \int _0^lv_t(x,t) \int _0^TP(x,t,\tau)u(x,\tau)d\tau

dxdt\biggr|\leq\\

\leq

 \int _0^T \int _0^lv_t^2dxdt+\frac{p}{(1-\sqrt{hT})^2} \int _0^T \int _0^lv^2dxdt.

\end{multline}

Теперь из \eqref{kir:8} имеем

$$

 \int _0^T \int _0^l\left((2b-K_t-1)v_t^2-a_tv_x^2-\Bigl(c_t+\frac{p}{(1-\sqrt{hT})^2}\Bigr)v^2\right)dxdt\leq

0,

$$

откуда в силу условий теоремы сразу следует, что $v(x,t)\equiv 0$

в $Q_T$.

\end{proof}



\smallskip



В силу доказанной эквивалентности задач 1 и 2 справедлива



\smallskip



\begin{theorem}[2]Если выполняются условия теоремы~$1,$ то существует не более одного

решения задачи~$1.$

\end{theorem}





\begin{thebibliography}{10}



\RBibitem{kir:sam}

\paper О некоторых проблемах современной теории дифференциальных

уравнений

\by А.~А.~Самарский

\jour Диффер. уравн.

\yr 1980

\vol 16

\issue 11

\pages 1925--1935

\misctransl 

\by A.~A.~Samarskiy

\paper  Some problems of the theory of differential equations

\jour Differ. Uravn.

\vol 16 

\yr 1980

\issue 11

\pages 1925–1935





\Bibitem{kir:cannon}

\paper The solution of the heat equation subject to the specification of energy

\by J.~R.~Cannon

\jour Quart. Appl. Math.

\yr 1963

\vol 21

\pages 155--160





\RBibitem{kir:Gu}

\by А.~К.~Гущин, В.~П.~Михайлов

\paper О разрешимости нелокальных задач для эллиптического уравнения второго порядка

\jour Матем. сб.

\yr 1994

\vol 185

\issue 1

\pages 121--160

\mathnet{http://mi.mathnet.ru/msb876}

\mathscinet{http://www.ams.org/mathscinet-getitem?mr=1264079}

\zmath{http://zbmath.org/?q=an:0827.35034}

\transl англ. пер.:~

\by 

\by A.~K.~Gushchin, V.~P.~Mikhailov

\paper On solvability of nonlocal problems for a~second-order elliptic equation

\jour Russian Acad. Sci. Sb. Math.

\yr 1995

\vol 81

\issue 1

\pages 101--136

\crossref{http://dx.doi.org/10.1070/SM1995v081n01ABEH003617}

\isi{http://gateway.isiknowledge.com/gateway/Gateway.cgi?GWVersion=2&SrcApp=PARTNER_APP&SrcAuth=LinksAMR&DestLinkType=FullRecord&DestApp=ALL_WOS&KeyUT=A1995QZ14400007}



\RBibitem{kir:skub}

\by А.~Л.~Скубачевский

\paper Неклассические краевые задачи.~I

\serial СМФН

\yr 2007

\vol 26

\pages 3--132

\publ РУДН

\publaddr М.

\mathnet{http://mi.mathnet.ru/cmfd116}

\mathscinet{http://www.ams.org/mathscinet-getitem?mr=2373390}

\zmath{http://zbmath.org/?q=an:1162.35022}

\transl англ. пер.:~

\by A.~L.~Skubachevskiy

\paper Nonclassical boundary value problems.~I. 

\jour Journal of Mathematical Sciences

\yr 2008

\vol 155

\issue 2

\pages 199--334

\crossref{http://dx.doi.org/10.1007/s10958-008-9218-9}



\RBibitem{kir:G-A}

\by Д.~Г.~Гордезиани, Г.~А.~Авалишвили

\paper Решения нелокальных задач для одномерных колебаний среды

\jour Матем. моделирование

\yr 2000

\vol 12

\issue 1

\pages 94--103

\mathnet{http://mi.mathnet.ru/mm832}

\mathscinet{http://www.ams.org/mathscinet-getitem?mr=1773208}

\zmath{http://zbmath.org/?q=an:1027.74505}

\misctransl

\by D.~G.~Gordeziani, G.~A.~Avalishvili

\paper On the constructing of solutions of the nonlocal initial boundary value problems for one-dimensional medium oscillation equations

\jour Matem. Mod.

\yr 2000

\vol 12

\issue 1

\pages 94--103



\RBibitem{kir:K-P}

\paper  О разрешимости краевых задач с~нелокальным граничным условием интегрального вида для многомерных гиперболических уравнений

\by А.~И.~Кожанов, Л.~С.~Пулькина

\jour Диффер. уравн.

\yr 2006

\vol 42

\issue 9

\pages 1166--1179

\transl англ. пер.:~

\by A.~I.~Kozhanov, L.~S.~Pul'kina

\paper On the solvability of boundary value problems with a nonlocal boundary condition of integral form for multidimensional hyperbolic equations

\jour Differ. Equ. 

\vol 42

\issue 9

\pages 1233--1246 

\yr 2006





\RBibitem{kir:Pu_1}

\by Л.~С.~Пулькина

\paper Краевые задачи для гиперболического уравнения с~нелокальными условиями~I и~II~рода

\jour Изв. вузов. Матем.

\yr 2012

\issue 4

\pages 74--83

\mathnet{http://mi.mathnet.ru/ivm8596}

\mathscinet{http://www.ams.org/mathscinet-getitem?mr=3076543}

\transl англ. пер.:~

\paper Boundary-value problems for a hyperbolic equation with nonlocal conditions of the I and II kind

\by L.~S.~Pul'kina 

\jour Russian Math. (Iz. VUZ)

\yr 2012

\vol 56

\issue 4

\pages 62--69

\crossref{http://dx.doi.org/10.3103/S1066369X12040081}





\RBibitem{kir:sabit}

\paper  Краевая задача для уравнения параболо-гиперболического типа с нелокальным интегральным условием

\by К.~Б.~Сабитов

\jour Диффер. уравн.

\yr 2010

\vol 46

\issue 10

\transl англ. пер.:~

\pages 1468--1478

\by K.~B.~Sabitov

\paper Boundary value problem for a parabolic-hyperbolic equation with a nonlocal integral condition

\jour  Differ. Equ. 

\vol 46

\issue 10

\pages 1472--1481 

\yr 2010





\end{thebibliography}





\noindent

\hskip6.5cm \thedate \par\noindent

\hskip6.5cm\therevdate





\makeengtitle



\newpage



%\end{document}

